% GENERATED FILE. Edit canonical lessons, then run npm run generate:books.
% canonical-source-hash: fnv1a64:f2f5be8033acf8c1
% canonical-lessons: UR-C133-javan, UR-C133-burha, UR-C133-pyara, UR-C133-narm, UR-C133-sakht

\chapter{Young, Old, Lovely, Soft, Hard}
\label{ch:ur-young-old-lovely-soft-hard}
\hlchaptermodality{133}

\begin{chapteropening}
\textbf{By the end of this chapter:} \emph{I can say young, old, lovely, soft and hard.}

Complete the last lesson of chapter 133: I can say young, old, lovely, soft and hard.
\end{chapteropening}

\section[\texorpdfstring{\ur{جوان} (javān)}{جوان (javān)}]{\ur{جوان} (javān) — young}

\label{lesson:UR-C133-javan}



\begin{quote}
Before the new one: say the Urdu for strong, then the Urdu for weak.
\end{quote}

\subsection*{You'll want to know: \ur{جوان}}
\textbf{\ur{جوان}} — \emph{javān} — \textquotedblleft{}young\textquotedblright{}.

\begin{cousinweb}
\textbf{\ur{جوان}} (\emph{javān}) is usually traced to Persian \emph{javān}, \textquotedblleft{}young\textquotedblright{}. CDIAL also lists a Hindi \emph{juwān}, \textquotedblleft{}a youth\textquotedblright{}, under \emph{*yuvāna}, \textquotedblleft{}young\textquotedblright{}, the Indo-Aryan form of the same ancient word (CDIAL 10510).
\end{cousinweb}

\begin{grammarlens}[title={What you've built}]
One more word for this chapter.
\end{grammarlens}

\subsection*{Your turn}
\begin{itemize}
\raggedright
  \item \emph{Say it:} \emph{javān}
  \item \emph{Say it:} \emph{javān}, once more
  \item \emph{Say it:} say \emph{javān}
  \item \emph{Recall:} say the Urdu for strong, then the Urdu for weak, then say \emph{javān} again
\end{itemize}

\begin{tcolorbox}[breakable,colback=teal!4,colframe=teal!35!black,title={Before you move on}]
What does \ur{جوان} mean? (\textquotedblleft{}Young\textquotedblright{}.) Say it once more.
\end{tcolorbox}
\section[\texorpdfstring{\ur{بوڑھا} (būṛhā)}{بوڑھا (būṛhā)}]{\ur{بوڑھا} (būṛhā) — old (of people)}

\label{lesson:UR-C133-burha}



\begin{quote}
Before the new one: say the Urdu for weak, then the Urdu for young.
\end{quote}

\subsection*{You'll want to know: \ur{بوڑھا}}
\textbf{\ur{بوڑھا}} — \emph{būṛhā} — \textquotedblleft{}old (of people)\textquotedblright{}.

\begin{cousinweb}
\textbf{\ur{بوڑھا}} (\emph{būṛhā}) is a tadbhava, a word inherited and reshaped by sound change. It comes from \emph{*buḍḍha}, \textquotedblleft{}old\textquotedblright{}, a form CDIAL reconstructs from its descendants (CDIAL 9271). CDIAL notes that the initial \emph{b-} rules out direct descent from Sanskrit \emph{vṛddha}, \textquotedblleft{}grown old\textquotedblright{}, though that word may have influenced it.
\end{cousinweb}

\begin{grammarlens}[title={What you've built}]
One more word for this chapter.
\end{grammarlens}

\subsection*{Your turn}
\begin{itemize}
\raggedright
  \item \emph{Say it:} \emph{būṛhā}
  \item \emph{Say it:} \emph{būṛhā}, once more
  \item \emph{Say it:} say \emph{būṛhā}
  \item \emph{Recall:} say the Urdu for weak, then the Urdu for young, then say \emph{būṛhā} again
\end{itemize}

\begin{tcolorbox}[breakable,colback=teal!4,colframe=teal!35!black,title={Before you move on}]
What does \ur{بوڑھا} mean? (\textquotedblleft{}Old (of people)\textquotedblright{}.) Say it once more.
\end{tcolorbox}
\section[\texorpdfstring{\ur{پیارا} (pyārā)}{پیارا (pyārā)}]{\ur{پیارا} (pyārā) — lovely, dear}

\label{lesson:UR-C133-pyara}



\begin{quote}
Before the new one: say the Urdu for young, then the Urdu for old.
\end{quote}

\subsection*{You'll want to know: \ur{پیارا}}
\textbf{\ur{پیارا}} — \emph{pyārā} — \textquotedblleft{}lovely, dear\textquotedblright{}.

\begin{cousinweb}
\textbf{\ur{پیارا}} (\emph{pyārā}) is a tadbhava, a word inherited from Sanskrit and reshaped by sound change. It comes from \emph{priyakāra}, \textquotedblleft{}doing a kindness\textquotedblright{} (CDIAL 8975). It is built from \emph{priya}, \textquotedblleft{}dear\textquotedblright{}, and \emph{kāra}, \textquotedblleft{}doing\textquotedblright{}.
\end{cousinweb}

\begin{grammarlens}[title={What you've built}]
One more word for this chapter.
\end{grammarlens}

\subsection*{Your turn}
\begin{itemize}
\raggedright
  \item \emph{Say it:} \emph{pyārā}
  \item \emph{Say it:} \emph{pyārā}, once more
  \item \emph{Say it:} say \emph{pyārā}
  \item \emph{Recall:} say the Urdu for young, then the Urdu for old, then say \emph{pyārā} again
\end{itemize}

\begin{tcolorbox}[breakable,colback=teal!4,colframe=teal!35!black,title={Before you move on}]
What does \ur{پیارا} mean? (\textquotedblleft{}Lovely, dear\textquotedblright{}.) Say it once more.
\end{tcolorbox}
\section[\texorpdfstring{\ur{نرم} (narm)}{نرم (narm)}]{\ur{نرم} (narm) — soft}

\label{lesson:UR-C133-narm}



\begin{quote}
Before the new one: say the Urdu for old, then the Urdu for lovely.
\end{quote}

\subsection*{You'll want to know: \ur{نرم}}
\textbf{\ur{نرم}} — \emph{narm} — \textquotedblleft{}soft\textquotedblright{}.

\begin{cousinweb}
\textbf{\ur{نرم}} (\emph{narm}) is a loan, usually traced to Persian \emph{narm}, \textquotedblleft{}soft\textquotedblright{}.
\end{cousinweb}

\begin{grammarlens}[title={What you've built}]
One more word for this chapter.
\end{grammarlens}

\subsection*{Your turn}
\begin{itemize}
\raggedright
  \item \emph{Say it:} \emph{narm}
  \item \emph{Say it:} \emph{narm}, once more
  \item \emph{Say it:} say \emph{narm}
  \item \emph{Recall:} say the Urdu for old, then the Urdu for lovely, then say \emph{narm} again
\end{itemize}

\begin{tcolorbox}[breakable,colback=teal!4,colframe=teal!35!black,title={Before you move on}]
What does \ur{نرم} mean? (\textquotedblleft{}Soft\textquotedblright{}.) Say it once more.
\end{tcolorbox}
\section[\texorpdfstring{\ur{سخت} (sakht)}{سخت (sakht)}]{\ur{سخت} (sakht) — hard}

\label{lesson:UR-C133-sakht}



\begin{quote}
Before the new one: say the Urdu for lovely, then the Urdu for soft.
\end{quote}

\subsection*{You'll want to know: \ur{سخت}}
\textbf{\ur{سخت}} — \emph{sakht} — \textquotedblleft{}hard\textquotedblright{}.

\begin{cousinweb}
\textbf{\ur{سخت}} (\emph{sakht}) is a loan, usually traced to Persian \emph{sakht}, \textquotedblleft{}hard, firm\textquotedblright{}.
\end{cousinweb}

\begin{grammarlens}[title={What you've built}]
One more word for this chapter.
\end{grammarlens}

\subsection*{Your turn}
\begin{itemize}
\raggedright
  \item \emph{Say it:} \emph{sakht}
  \item \emph{Say it:} \emph{sakht}, once more
  \item \emph{Say it:} say \emph{sakht}
  \item \emph{Recall:} say the Urdu for lovely, then the Urdu for soft, then say \emph{sakht} again
\end{itemize}

\begin{tcolorbox}[breakable,colback=teal!4,colframe=teal!35!black,title={Before you move on}]
What does \ur{سخت} mean? (\textquotedblleft{}Hard\textquotedblright{}.) Say it once more.
\end{tcolorbox}
